\chapter{SAM and OAM interaction (SOI) of LG Beam}
\label{ch:Theory}
% 連結SHG&DFT
% 2-1
\section{Gauge \add{freedom}}

\add{Although electromagnetic fields are the directly observable physical quantities, structured light beams such as Laguerre--Gaussian (LG) beams are often formulated in terms of the scalar and vector potentials. However, the potentials possess gauge freedom, leading to different forms of the wave equations under different gauge conditions. Importantly, the derivation of LG beams relies on the paraxial approximation, whose implementation is not completely independent of the chosen gauge.}

In this section, we examine how gauge choice \add{affects} the formulation of wave equations. We begin by introducing the wave equations under different gauge conditions, and then analyze the compatibility of the paraxial approximation. This discussion clarifies why certain gauges are more suitable than others for describing paraxial optical fields.

\subsection{Gauge Freedom in Wave Equations}

Maxwell's equations describe how the \textbf{electric field} $\boldsymbol{E}(\boldsymbol{r}, t)$ and the \textbf{magnetic field} $\boldsymbol{B}(\boldsymbol{r}, t)$ evolve in space and time. These equations form the fundamental basis for describing electromagnetic waves, including Laguerre--Gaussian beams. In \textbf{free space}, the charge density $\rho$ and current density $\mathbf{J}$ vanish. In this case, Maxwell's equations reduce to four homogeneous differential equations involving the vacuum permittivity $\epsilon_0$, the vacuum permeability $\mu_0$, and the speed of light in vacuum $c=1/\sqrt{\mu_0\epsilon_0}$. 

Within this electrodynamics framework, Maxwell's equations can be rewritten as two coupled wave equations for the \textbf{scalar and vector potentials}, $\phi(\boldsymbol{r}, t)$ and $\boldsymbol{A}(\boldsymbol{r}, t)$, derived in appendix~\ref{ch:Maxwell} as
\begin{align} 
&\nabla^2\phi(\boldsymbol{r}, t)+\frac{\partial }{\partial t}\left(\boldsymbol{\nabla}\cdot\boldsymbol{A}(\boldsymbol{r}, t)\right)=0, \label{max5} \\
&\nabla^2\boldsymbol{A}(\boldsymbol{r}, t)-\frac{1}{c^2}\frac{\partial^2 }{\partial t^2}\boldsymbol{A}(\boldsymbol{r}, t)-\boldsymbol{\nabla}\left(\boldsymbol{\nabla}\cdot\boldsymbol{A}(\boldsymbol{r}, t)+\frac{1}{c^2}\frac{\partial }{\partial t}\phi(\boldsymbol{r}, t)\right)=0. \label{max6}
\end{align}
These equations contain the same physical information as the original Maxwell equations, but are expressed in terms of the scalar and vector potentials.

Further simplification of Eqs.~(\ref{max5}) and (\ref{max6}) can be achieved by exploiting the gauge freedom of the potentials. Below, we introduce two commonly used gauges: the \textbf{Coulomb gauge} and the \textbf{Lorenz gauge}. Superscripts $C$ and $L$ are used to distinguish quantities defined in the Coulomb and Lorenz gauges, respectively. The gauge conditions are given by
\begin{align} 
\text{Coulomb gauge condition}: \quad &\boldsymbol{\nabla}\cdot\boldsymbol{A}^C(\boldsymbol{r}, t)=0,  \label{Ccondition}\\
\text{Lorenz gauge condition}: \quad &\boldsymbol{\nabla}\cdot\boldsymbol{A}^L(\boldsymbol{r}, t)=-\frac{1}{c^2}\frac{\partial }{\partial t}\phi^L(\boldsymbol{r}, t).
\end{align}
Both gauges simplify the wave equations. In the Coulomb gauge, they reduce to
\begin{align} 
&\nabla^2\phi^C(\boldsymbol{r}, t)=0, \text{and} \label{Cwave1}\\
&\nabla^2\boldsymbol{A}^C(\boldsymbol{r}, t)-\frac{1}{c^2}\frac{\partial^2 }{\partial t^2}\boldsymbol{A}^C(\boldsymbol{r}, t)=\frac{1}{c^2}\frac{\partial }{\partial t}\boldsymbol{\nabla}\phi^C(\boldsymbol{r}, t)=0, \label{Cwave2}
\end{align}
where the last equality follows from Eq.~(\ref{Cwave1}).
In the Lorenz gauge, Eqs.~(\ref{max5}) and (\ref{max6}) become
\begin{align} 
&\nabla^2\phi^L(\boldsymbol{r}, t)-\frac{1}{c^2}\frac{\partial^2 }{\partial t^2}\phi^L(\boldsymbol{r}, t)=0, \text{and} \label{Lwave1}\\
&\nabla^2\boldsymbol{A}^L(\boldsymbol{r}, t)-\frac{1}{c^2}\frac{\partial^2 }{\partial t^2}\boldsymbol{A}^L(\boldsymbol{r}, t)=0. \label{Lwave2}
\end{align}
\add{Here, we introduce the \textbf{monochromatic wave} ansatz, in which the electromagnetic fields and potentials are expressed as 
$\boldsymbol{E}(\boldsymbol{r}, t)=\mathrm{Re}\left[\boldsymbol{E}(\boldsymbol{r}) e^{-i\omega t}\right]$,
$\boldsymbol{A}^C(\boldsymbol{r}, t)=\mathrm{Re}\left[\boldsymbol{A}^C(\boldsymbol{r}) e^{-i\omega t}\right]$,
$\boldsymbol{A}^L(\boldsymbol{r}, t)=\mathrm{Re}\left[\boldsymbol{A}^L(\boldsymbol{r}) e^{-i\omega t}\right]$, and
$\phi^L(\boldsymbol{r}, t)=\mathrm{Re}\left[\phi^L(\boldsymbol{r}) e^{-i\omega t}\right]$,
where $\omega$ is the angular frequency of light. Substituting these expressions into Eq.~(\ref{gaugeF}) from Maxwell's equations and applying the corresponding gauge conditions leads to}
\begin{equation} 
\boldsymbol{E}(\boldsymbol{r})=i\omega\boldsymbol{A}^C(\boldsymbol{r})=
i\omega\left(\boldsymbol{A}^L(\boldsymbol{r})+
\frac{\boldsymbol{\nabla}\left(\boldsymbol{\nabla}\cdot\boldsymbol{A}^L(\boldsymbol{r})\right)}{q_0^2}\right).
\label{LCrelation}
\end{equation}
\add{Here $q_0=\omega/c$ denotes the wavenumber of light in free space. Under the monochromatic-wave ansatz, electric field $\boldsymbol{E}(\boldsymbol{r}, t)$ and magnetic field $\boldsymbol{B}(\boldsymbol{r}, t)$ are related to potentials through Eq.~(\ref{LCrelation}) and Eq.~(\ref{BcrulA})}
%
%
%
%
%
%
\subsection{Incompatibility of Gauge Freedom with the Paraxial Approximation}
In principle, the vector potential formulation of LG beams can be derived starting from either Eq.~(\ref{Cwave2}) or Eq.~(\ref{Lwave2}). \add{However, for a transversely polarized wave, the directional constraint imposed by the Coulomb gauge is ultimately incompatible with the standard paraxial Helmholtz formulation. For clarity, this incompatibility will be discussed in detail after deriving the paraxial Helmholtz equation within the Lorenz gauge framework~\cite{andrews2012angular}.}

\add{The derivation starts from the scalar Helmholtz equation, which is obtained from} Eq.~(\ref{Lwave2}) by imposing both the monochromatic ($\boldsymbol{A}^L(\boldsymbol{r}, t)=\text{Re}\left[\boldsymbol{A}^L(\boldsymbol{r}) e^{-i\omega t}\right]$) and transverse wave ans\"atze ($\boldsymbol{A}^L(\boldsymbol{r})=\hat{\boldsymbol{\varepsilon}}A^L(\boldsymbol{r})$, where we condsider polarization unit vector $\hat{\boldsymbol{\varepsilon}}=\hat{\boldsymbol{\varepsilon}}_{\sigma}=\hat{\boldsymbol{x}}+i \sigma \hat{\boldsymbol{y}},\,\sigma=\pm 1$ as right/left circularly polarized)(see Sec.~\ref{ch:SHE}):

\begin{equation}
\nabla^2 A^L(\boldsymbol{r}) + q_0^2 A^L(\boldsymbol{r}) = 0.
\label{SH}
\end{equation}
The \textbf{paraxial approximation}~\cite{andrews2012angular, kasap2013optoelectronics, pedrotti2007introduction} assumes that the rays have only small inclinations with respect to the optical axis in geometric optics. In wave optics, this corresponds to an angular spectrum dominated by plane-wave components whose wavevectors are close to the propagation direction. Here, we choose the propagation direction along the $\boldsymbol{z}$ axis, such that the dominant component of the wavevector $\boldsymbol{q}_0$ lies along $\hat{\boldsymbol{z}}$.

\begin{wrapfigure}{l}{0.325\textwidth}
    \centering
    \includegraphics[width=\linewidth, trim=8.3cm 5.8cm 8.3cm 5.8cm, clip]{wave_vector.pdf} %左 下 右 上 
    \caption{\add{Illustration of a dominant component of wavevector $\boldsymbol{q}_0$ lies along $\hat{\boldsymbol{z}}$}}
    \label{fig:wave_vector}
\end{wrapfigure}

Under this approximation, we introduce the \textbf{paraxial wave} ansatz, $A^L(\boldsymbol{r}) = U(\boldsymbol{r}) e^{iq_0 z}$, and substitute it into Eq.~(\ref{SH}), which yields 
\begin{equation}
\nabla^2_\parallel U(\boldsymbol{r})
+ \frac{\partial^2}{\partial z^2} U(\boldsymbol{r})
+ 2 i q_0 \frac{\partial}{\partial z} U(\boldsymbol{r}) = 0.
\end{equation}
Here, $\nabla^2_\parallel = \frac{\partial^2}{\partial x^2} + \frac{\partial^2}{\partial y^2}$ denotes the transverse Laplacian. The paraxial wave ansatz assumes that the \textbf{envelope function} $U(\boldsymbol{r})$ \textbf{varies slowly along the} $\boldsymbol{z}$ \textbf{direction} compared to the rapid phase variation of $e^{i q_0 z}$. Consequently, the paraxial approximation implies the following relations:

\begin{minipage}{0.48\textwidth}
    \begin{equation}
    \left| \frac{\partial^2}{\partial z^2} U(\boldsymbol{r}) \right|
    \ll
    \left| \nabla^2_\parallel U(\boldsymbol{r}) \right|,
    \label{paraxial1}
    \end{equation}
\end{minipage}
\hfill % 自動填滿中間的空白，讓兩個公式漂亮地分開
\begin{minipage}{0.48\textwidth}
    \begin{equation}
    \left| \frac{\partial^2}{\partial z^2} U(\boldsymbol{r}) \right|
    \ll
    q_0 \left| \frac{\partial}{\partial z} U(\boldsymbol{r}) \right|.
    \label{paraxial2}
    \end{equation}
\end{minipage}
\vspace{1.75\baselineskip} % 給下方內文留一點標準行高空隙
Neglecting the second-order derivative with respect to $z$ then leads to the paraxial Helmholtz equation
\begin{equation}
\nabla^2_\parallel U(\boldsymbol{r}) + 2 i q_0 \frac{\partial}{\partial z} U(\boldsymbol{r}) \color{G}{ \approx} 0.
\label{PHE}
\end{equation}

\add{We now clarify why a transversely polarized light ($\hat{\boldsymbol{\varepsilon}}$ confined to the $x$--$y$ plane) is not suitable in this derivation under the Coulomb gauge. For a transversely polarized wave, the Coulomb gauge condition gives}
\begin{align} 
\color{G} 
\boldsymbol{\nabla} \cdot \boldsymbol{A}^C(\boldsymbol{r})
&=
\hat{\boldsymbol{\varepsilon}}_{\sigma} \cdot \boldsymbol{\nabla}_\parallel A^C(\boldsymbol{r}) \\
&=
\left( \frac{\partial}{\partial x}\hat{\boldsymbol{x}} + i \sigma \frac{\partial}{\partial y}\hat{\boldsymbol{y}} \right)A^C(\boldsymbol{r}) = 0
\label{Vio}
\end{align}
\add{where $\boldsymbol{\nabla}_\parallel = \frac{\partial}{\partial x}\hat{\boldsymbol{x}} + \frac{\partial}{\partial y}\hat{\boldsymbol{y}}$. This condition imposes a directional constraint in the transverse plane on the vector potential and results in the transverse Laplaian of the potential become zero.}
\begin{equation}
\color{G}
\nabla^2_\parallel A^C(\boldsymbol{r})
=
\left( \frac{\partial}{\partial x} + i \frac{\partial}{\partial y} \right)
\left( \frac{\partial}{\partial x} - i \frac{\partial}{\partial y} \right)
A^C(\boldsymbol{r})
= 0,
\end{equation}
\add{which in turn is incompatible with the approximation in Eq.~(\ref{paraxial1})~\cite{andrews2012angular}. In contrast, Lorenz gauge do not  have such constraint in transverse plane. To satisfy the divergence-free condition in the Coulomb gauge, a nonzero longitudinal ($z$) component of the vector potential should be introduced, going beyond the simple transverse polarization. Therefore, the Lorenz gauge provides a more natural and consistent framework for deriving the paraxial Helmholtz equation and Laguerre--Gaussian beam solutions.}
%
%
%
%
%
%
\section{Manifestation of SOI in LG Beam}

Within the paraxial approximation, the optical helicity $\sigma$ and the topological charge $\ell$ are directly proportional to the spin and orbital angular momentum of light, respectively~\cite{bliokh2015transverse, barnett1994orbital}. This correspondence provides a convenient framework for analyzing the angular momentum properties of LG beams. 

Moreover, when the vector potential formulation of LG beams are expressed in the Coulomb gauge, an SOI can emerge through the additional longitudinal component of the field, giving rise to terms proportional to $\sigma \ell$ and $\sigma + \ell$. Adopting the Coulomb gauge is further justified by the theory of light–matter interaction, where the vector potential is conventionally formulated within both the Coulomb gauge and the angular spectrum representation. \add{Where in angular spectrum representation can we relate the in-plane wavevector to the vector in reciprocal space, which connect to theory of matter.}

In the following section, we first introduce the LG beams solutions in both real-space and angular spectrum representations. We then perform a gauge transformation to reveal the SOI.
%%%%%%%%%%
\begin{comment} 
can be solved in various coordinate systems and yields different beam solutions, such as the \textbf{Hermite--Gaussian} beam, the \textbf{Laguerre--Gaussian} (LG) beam, and the \textbf{Bessel} beam. 
In cylindrical coordinates $(\rho,\phi,z)$, the paraxial Helmholtz equation takes the form
\begin{equation}
\left(
\frac{1}{\rho}\frac{\partial}{\partial\rho}
+
\frac{\partial^2}{\partial\rho^2}
+
\frac{1}{\rho^2}\frac{\partial^2}{\partial\phi^2}
+
2iq_0\frac{\partial}{\partial z}
\right)
U(\rho,\phi,z)=0 .
\end{equation}
\end{comment}
%%%%%%%%%%

\subsection{Solutions of LG Beams}
Under the \add{\textbf{long Rayleigh-range} approximation~\cite{romero2002quantum}}\footnote{\add{For a light beam, the Rayleigh range is the propagation distance from the beam waist to the point where the beam cross-sectional area doubles (see Figure~\ref{fig:Rrange}). It is given by $z_R=\frac{1}{2}q_0W_0^2$, where $W_0$ is the beam waist. When the Rayleigh range is much larger than the characteristic propagation distance of interest, diffraction-induced beam spreading becomes negligible. In this limit, the beam waist can be treated as approximately constant, and the envelope function may be regarded as independent of $z$. This simplification is referred to as the long Rayleigh-range approximation.}}, the LG beam solutions of the paraxial Helmholtz equation, Eq.~(\ref{PHE}), can be expressed in cylindrical coordinates as
\begin{equation}
\color{G}
\boldsymbol{A}^{L}_{\sigma, \ell, p}(\boldsymbol{r})
=
\hat{\boldsymbol{\varepsilon}}_\sigma  A^{L}_{\ell, p}(\boldsymbol{r})
=
\hat{\boldsymbol{\varepsilon}}_\sigma 
\, u_{\ell, p}(\boldsymbol{r}_\parallel)
\, e^{i q_0 z},
\label{LGbeam}
\end{equation}
where the position vector is written as $\boldsymbol{r}=(\rho,\phi,z)=(\boldsymbol{r}_\parallel,z)$. Here $\boldsymbol{r}_\parallel$ denotes the in-plane position vector consisting of the radial coordinate $\rho$ and the azimuthal angle $\phi$. The integer $\ell\in\mathbb{Z}$ is the \textbf{topological charge}, and the integer $p\ge0$ is the radial mode index. The vector $\boldsymbol{q}_0 = q_0 \hat{\boldsymbol{z}}$ denotes the wavevector along the propagation direction. The \add{envelope} function is expressed as
\begin{equation}
\color{G}
u_{\ell, p}(\boldsymbol{r}_\parallel)=f_{\ell, p}(\rho)\,e^{i\ell\phi},
\end{equation}
where the radial function is $f_{\ell, p}(\rho)
=
A^{\rm LG}_0\sqrt{\frac{2p!}{\pi(|\ell|+p)!}}
\,
L_p^{|\ell|}
\left(
\frac{2\rho^2}{W_0^2}
\right)
\left(
\frac{\sqrt{2}\rho}{W_0}
\right)^{|\ell|}
\exp\left(
-\frac{\rho^2}{W_0^2}
\right)$, which is given by~\cite{romero2002quantum, peng2022tailoring, sanchez2024enhanced} and $A^{\rm LG}_0$ is a constant amplitude. Here the generalized Laguerre polynomial is defined as $L_p^{|\ell|}(x)
=
\sum_{m=0}^{p}
(-1)^m
\frac{(|\ell|+p)!}{(p-m)!(|\ell|+m)!m!}
x^m$, with $x=\frac{2\rho^2}{W_0^2}$. The parameter $W_0$ denotes the beam waist of the LG beam, while the cylindrical coordinates are related to the Cartesian coordinates through $\rho=\sqrt{x^2+y^2}$ and $\phi=\tan^{-1}(y/x)$. \add{The geometric meaning of the beam waist and the Rayleigh range is illustrated in Figure~\ref{fig:Rrange}.}

\begin{wrapfigure}{l}{0.455\textwidth}
\centering
\includegraphics[width=\linewidth, trim=1.5cm 3cm 1.5cm 3cm, clip]{Rrange.pdf} %左 下 右 上
\caption{\add{Illustration of the evolution of the beam radius $W(z)$ during propagation. The minimum beam radius $W_0$ occurs at the beam waist, and the Rayleigh range $z_R$ is defined as the distance from the waist at which $W(z_R)=\sqrt{2}W_0$.}}
\label{fig:Rrange}
\end{wrapfigure}

Focusing on the minimum cross-section area of the beam ($z=0$), Figure~\ref{fig:ufig} shows that as $\ell$ and/or $p$ increase, $u_{\ell, p}(\boldsymbol{r}_\parallel)$ develops a broader doughnut-shaped intensity profile, accompanied by an increasing number of radial nodes along the radial direction. The phase varies continuously with the azimuthal angle $\phi$, reflecting the helical wavefront associated with orbital angular momentum. In particular, the phase winds by $2\pi \ell$ around the beam center. The beam waist is chosen as $W_0 = 1.5~\mu\mathrm{m}$, which sets the spatial scale of the beam and determines the radial extent of the intensity distribution. \\

\begin{figure}[H]
\centering
\includegraphics[width=0.95\linewidth, trim=0.2cm 0cm 0.2cm 0.2cm, clip]{ufig.pdf}%左 下 右 上
\caption{Squared magnitude (intensity) and real part distribution of $A^{L}_{\ell, p}(\boldsymbol{r})$ at $z=0$. The green dotted line denotes the beam waist $W_0$. The scale bar indicates the length of 1 $\mathrm{\mu m}$.}
\label{fig:ufig}
\end{figure}

%
%
%
%
%
%
%
%

To derive the vector potential of LG beams in the angular spectrum representation, the field is expanded \textbf{in a plane-wave basis} as
\begin{equation}
\boldsymbol{A}(\boldsymbol{r}) = \sum_{\boldsymbol{q}} \boldsymbol{\mathcal{A}}(\boldsymbol{q}) e^{i \boldsymbol{q} \cdot \boldsymbol{r}}, 
\label{planwave}
\end{equation}
where $\boldsymbol{\mathcal{A}}(\boldsymbol{q})$ is the complex amplitude associated with each plane-wave component $e^{i\boldsymbol{q}\cdot\boldsymbol{r}}$. The LG beams in Eq.~(\ref{LGbeam}) can be written in this form through a defined \textbf{two-dimensional Fourier transformation} and its corresponding inverse transformation,
\begin{equation}
u_{\ell, p}(\boldsymbol{r}_\parallel) = \frac{1}{(2\pi)^2} \int {\rm d}^2 \boldsymbol{q}_\parallel\, 
F^{\rm LG}_{\ell, p}(\boldsymbol{q}_\parallel)\, e^{i\boldsymbol{q}_\parallel\cdot \boldsymbol{r}_\parallel},
\quad \longleftrightarrow \quad
F^{\rm LG}_{\ell, p}(\boldsymbol{q}_\parallel) = \int {\rm d}^2 \boldsymbol{r}_\parallel\, 
u_{\ell, p}(\boldsymbol{r}_\parallel)\, e^{-i \boldsymbol{q}_\parallel\cdot \boldsymbol{r}_\parallel}.
\label{eq:FourierPair}
\end{equation}
For applications to two-dimensional materials, the in-plane wavevector $\boldsymbol{q}_\parallel$ can be approximately identified with a reciprocal-space vector. Under this approximation, Eq.~(\ref{eq:FourierPair}) can be discretized, and Eq.~(\ref{planwave}) becomes
\begin{equation} 
\color{G}
\boldsymbol{\mathcal{A}}^{L}_{\sigma, \ell, p}(\boldsymbol{r})
=
\sum_{\boldsymbol{q}_\parallel}
\boldsymbol{\mathcal{A}}^L_{\sigma, \ell, p}(\boldsymbol{q}_\parallel)
e^{i\boldsymbol{q}_\parallel\cdot \boldsymbol{r}_\parallel}
e^{i q_0 z}.
\end{equation}
Here the in-plane wavevector is written as $\boldsymbol{q}_\parallel=(q_\parallel,\phi_{\boldsymbol{q}_\parallel})$. The corresponding complex amplitude can be obtained analytically, as summarized in Table~\ref{tab:my_formulas}.
\begin{table}[H]
\color{G}
\centering
\caption{Real space and angular spectrum representations of LG beams}
\label{tab:my_formulas}

% 縮減回 2 欄結構，確保橫向空間最小化
\begin{tabular}{|c|c|}
\hline
\textbf{Real space} & \textbf{Angular spectrum} \\ \hline

% 第一類別分隔列 
\hline
\multicolumn{2}{|c|}{\textbf{Vector potential}} \\ \hline
\rule[-2.5ex]{0pt}{6.5ex} % 調整支柱高度以配合 \displaystyle 的飽滿尺寸
$\displaystyle \boldsymbol{A}^{L}_{\sigma, \ell, p}(\boldsymbol{r})
=
\hat{\boldsymbol{\varepsilon}}_\sigma A^{L}_{\ell, p}(\boldsymbol{r})
=
\hat{\boldsymbol{\varepsilon}}_\sigma
\, u_{\ell, p}(\boldsymbol{r}_\parallel)
\, e^{i q_0 z}$
&
$\displaystyle \boldsymbol{\mathcal{A}}^L_{\sigma, \ell, p}(\boldsymbol{q}_\parallel)
=
\hat{\boldsymbol{\varepsilon}}_\sigma \mathcal{A}^L_{\ell, p}(\boldsymbol{q}_\parallel)
=
\hat{\boldsymbol{\varepsilon}}_\sigma
\frac{1}{\Omega}
\, F_{|\ell|, p}^{\rm LG}(\boldsymbol{q}_\parallel)$ \\ \hline

% 第二類別分隔列
\hline
\multicolumn{2}{|c|}{\textbf{Envelope function}} \\ \hline
\rule[-2.5ex]{0pt}{6.5ex}
$\displaystyle u_{\ell, p}(\boldsymbol{r}_\parallel)
=
f_{\ell, p}(\rho)e^{i \ell \phi}$
&
$\displaystyle \frac{1}{\Omega}F_{\ell, p}^{\rm LG}(\boldsymbol{q}_\parallel)
=
\tilde{F}_{\left|\ell\right|, p}(q_\parallel)e^{i \ell \phi_{q_\parallel}}$ \\ \hline

% 第三類別分隔列 (使用 \displaystyle 強制將換行公式內的字體放大至完全一致)
\hline
\multicolumn{2}{|c|}{\textbf{Solutions at $p=0$}} \\ \hline
\rule[-6ex]{0pt}{13ex} % 微調這列雙行公式的上下留白
$\begin{aligned} 
&\displaystyle f_{\ell, p=0}(\rho) = \\
&\displaystyle \frac{A^{\rm LG}_0}{\sqrt{(|\ell|)!\pi/2}}
\left( \frac{\sqrt{2}\rho}{W_0} \right)^{|\ell|} \exp\left( -\frac{\rho^2}{W_0^2} \right)
\end{aligned}$
&
$\begin{aligned} 
&\displaystyle \tilde{F}_{\left|\ell\right|, p=0}(q_\parallel) = \\
& \displaystyle \frac{A^{\rm LG}_0 W_0^2 \left(-i\right)^{\left|\ell\right|}\pi}{\Omega\sqrt{(|\ell|)!\pi/2}}
 \left( \frac{q_{\parallel} W_0}{\sqrt{2}} \right)^{|\ell|} \exp\left( -\frac{q_{\parallel}^2 W_0^2}{4} \right)
\end{aligned}$ \\ \hline 

\end{tabular}
\end{table}
\noindent
where
$
\tilde{F}_{\left|\ell\right|, p}(q_\parallel)
=
\frac{(-i)^{|\ell|}}{\Omega}
2\pi\,
\mathcal{H}_{|\ell|}\!\left[f_{\ell, p}(\rho)\right],
$
and
$
\mathcal{H}_{|\ell|}\!\left[f_{\ell, p}(\rho)\right]
=
\int_0^\infty {\rm d}\rho\;
\rho\, f_{\ell, p}(\rho)\,
J_{|\ell|}(q_\parallel\rho)
$
is the \textbf{Hankel transform} of order $|\ell|$, where $J_{|\ell|}$ is the \textbf{Bessel function of the first kind}. Here $\Omega$ denotes the area of the two-dimensional system

Figure~\ref{fig:ffig} shows that \add{$\mathcal{A}^L_{\ell, p}(\boldsymbol{q}_\parallel)$} exhibits an intensity profile analogous to that of \add{$A^{L}_{\ell, p}(\boldsymbol{r})$}, reflecting the characteristic doughnut-shaped structure in angular spectrum. This correspondence follows directly from the Fourier transform relation between two representations of the LG beam. The doughnut-shaped distribution in angular spectrum indicates that the dominant in-plane wavevectors lie within a finite range determined by the beam waist $W_0$.
\begin{figure}[H]
\centering
\includegraphics[width=1.0\linewidth]{ffig.pdf}
\caption{Squared magnitude (intensity) and real part distribution of $\mathcal{A}^L_{\ell, p=0}(\boldsymbol{q}_\parallel)$. The scale bar indicates the considered $0.1$ times light cone range for WSe$_2$-monolayers, $Q_c = 8.6~\mu\mathrm{m}^{-1}$.}
\label{fig:ffig}
\end{figure}

%
%
%
%
%
%
% 

\subsection{\add{Longitudinal Component of LG Beam : SOI effect}}

In the previous section, Eq.~(\ref{LCrelation}) was introduced to relate the vector potentials defined in the two gauges to the physical fields. Using this relation, we can perform gauge transformation by rewritten Eq.~(\ref{LCrelation}) as
\begin{equation}
\color{G}
\boldsymbol{A}^{C}_{\sigma, \ell, p}(\boldsymbol{r})
=
\boldsymbol{A}^{L}_{\sigma, \ell, p}(\boldsymbol{r})
+
\frac{\boldsymbol{\nabla}
\left(
\boldsymbol{\nabla}\cdot
\boldsymbol{A}^{L}_{\sigma, \ell, p}(\boldsymbol{r})
\right)}
{q_0^2}.
\label{LCrelation2}
\end{equation}
\add{To explicitly incorporate the SAM of light, the polarization unit vector is defined as
$\hat{\boldsymbol{\varepsilon}}\equiv\hat{\boldsymbol{\varepsilon}}_{\sigma}=\frac{1}{\sqrt{2}}(\hat{\boldsymbol{x}}+i\sigma\hat{\boldsymbol{y}})$,
where $\sigma=\pm1$ denotes the spin of the photon. From now on we omit the superscript $C$ to simplify the notations and the vector potential in real space under the Coulomb gauge is obtained as}
\begin{equation}
\color{G}
\boldsymbol{A}_{\sigma, \ell, p}(\boldsymbol{r})
\equiv 
\boldsymbol{A}^{C}_{\sigma, \ell, p}(\boldsymbol{r})
=\tilde{A}^{\parallel}_{\ell, p}(\boldsymbol{r}) \hat{\boldsymbol{\varepsilon}} + \tilde{A}^{z}_{\sigma, \ell, p}(\boldsymbol{r}) \hat{\boldsymbol{z}}.
\end{equation}
\add{For LG beams in the fundamental radial mode $p=0$} 
\begin{equation} 
\color{G}
\begin{aligned}
\tilde{A}^{\parallel}_{\ell, p}(\boldsymbol{r})
&\approx 
f_{\ell,p=0}(\rho)e^{i\ell \phi}e^{i q_0 z}\\
\tilde{A}^{z}_{\sigma, \ell, p}(\boldsymbol{r})
&\approx
\frac{i}{\sqrt{2}q_0\rho}
\left(
\bigl(|\ell|+\sigma \ell\bigr)-\frac{2\rho^2}{W_0^2}
\right)
f_{\ell,p=0}(\rho)
e^{i(\sigma+\ell)\phi}e^{i q_0 z}
\end{aligned}
\label{AinRinC} % \label 移到這裡，緊接在 \end{aligned} 後面
\end{equation}
In general, the second term in Eq.~(\ref{LCrelation2}) introduces corrections to all spatial components of the vector potential. However, under the paraxial approximation, these corrections are suppressed. Therefore, we retain only the longitudinal ($z$) component as the leading correction.

The longitudinal component in Eq.~(\ref{AinRinC}) reveals two distinct types of SOI: a $\sigma\ell$ dependence in the amplitude and a $\sigma+\ell$ dependence in the phase. Figure~\ref{fig:AinRinC} shows that the magnitude (absolute value) at $z=0$ of this longitudinal component is significantly weaker than that of the transverse components. The $\sigma\ell$ dependence manifests as an asymmetry between cases with the same $\ell$ but opposite spin $\sigma$.

It is worth noting that the longitudinal term is weighted by a factor proportional to $(q_0 W_0^2)^{-1}$, which is inversely related to the Rayleigh range ($\frac{1}{2} q_0 W_0^2$)~\cite{kasap2013optoelectronics} . This indicates that a smaller beam waist $W_0$ leads to a stronger longitudinal component, corresponding to a shorter Rayleigh range and tighter focusing, as shown in the bottom row of Figure~\ref{fig:AinRinC}. The beam waist $W_0$ is chosen to be $1.5~\mu\mathrm{m}$ and $0.5~\mu\mathrm{m}$, representing a typical paraxial beam and a case at the boundary of the paraxial regime, respectively~\cite{forbes2021relevance, klimov2012mapping, nemoto1990nonparaxial}. This demonstrates that the longitudinal component, although small, remains finite even within the paraxial regime.
\begin{figure}[H]
\centering
\includegraphics[width=1.0\linewidth]{AinRinC.pdf}
\caption{Magnitude (absolute value) of \add{$\boldsymbol{A}_{\sigma, \ell, p=0}(\boldsymbol{r})$} at $z=0$ for different components and spin. The $\sigma\ell$ term leads to asymmetry between cases with the same $\ell$ in the longitudinal component (a wavelength of $532~\mathrm{nm}$ is assumed).}
\label{fig:AinRinC}
\end{figure}

%
%
%
%
As for the $\sigma+\ell$ dependence in the phase term, it is retained in the angular spectrum representation and will be discussed below. By combining Eq.~(\ref{LCrelation2}) with the two-dimensional Fourier transformation defined in Coulomb gauge, vector potential can be written as
\begin{equation}
\color{G}
\boldsymbol{A}_{\sigma, \ell, p}(\boldsymbol{r})
=
\frac{1}{(2\pi)^2}
\int {\rm d}^2 \boldsymbol{q}_\parallel
\,
\boldsymbol{\mathcal{A}}_{\sigma, \ell, p}(\boldsymbol{q_\parallel})
e^{i\boldsymbol{q_\parallel}\cdot \boldsymbol{r}_\parallel}
e^{iq_0 z}.
\label{planwaveinC}
\end{equation}
\add{Since we primarily focus on the fundamental radial mode $p=0$, from now on we omit the subscript $p$ to simplify the notations. The complex amplitude under Coulomb gauge in the angular spectrum representation is then written as
\begin{equation}
\color{G}
\boldsymbol{\mathcal{A}}_{\sigma, \ell}(\boldsymbol{q}_\parallel)
=\boldsymbol{\mathcal{A}}_{\sigma, \ell, p=0}(\boldsymbol{q}_\parallel)
=\tilde{\mathcal{A}}^{\parallel}_{\ell}(\boldsymbol{q}_\parallel) \hat{\boldsymbol{\varepsilon}} + \tilde{\mathcal{A}}^{z}_{\sigma, \ell}(\boldsymbol{q}_\parallel) \hat{\boldsymbol{z}}.
\label{AinCinQ}
\end{equation}}
The corresponding complex amplitude in the angular spectrum representation is obtained as
\begin{align}
\color{G}
\boldsymbol{\mathcal{A}}_{\sigma, \ell, p}(\boldsymbol{q_\parallel})
&=
\left[
\hat{\boldsymbol{\varepsilon}}
-
\left(
\frac{|\boldsymbol{q}|}{q_0}
\right)^2
(\hat{\boldsymbol{\varepsilon}}\cdot\hat{\boldsymbol{q}_\parallel})
\hat{\boldsymbol{q}}
\right]
\mathcal{A}^L_{\ell, p}(\boldsymbol{q_\parallel}) \\
&=
\left[
\hat{\boldsymbol{\varepsilon}}
-
\left(
\frac{|\boldsymbol{q}|}{q_0}
\right)^2
\frac{\sin\theta_{\boldsymbol{q}}}{\sqrt{2}}
e^{i\sigma\phi_{\boldsymbol{q}_\parallel}}
\hat{\boldsymbol{q}}
\right]
\mathcal{A}^L_{\ell, p}(\boldsymbol{q_\parallel}) \label{sigma_limit} \\
&\approx
\left[
\hat{\boldsymbol{\varepsilon}}
-
\frac{\sin\theta_{\boldsymbol{q}}}{\sqrt{2}}
e^{i\sigma\phi_{\boldsymbol{q}_\parallel}}
\hat{\boldsymbol{z}}
\right]
\mathcal{A}^L_{\ell, p}(\boldsymbol{q_\parallel}),
\label{approx_to_1}
\end{align}
where the final approximation follows from the paraxial condition $\cos\theta_{\boldsymbol{q}} = q_0/|\boldsymbol{q}| \approx 1$. The wavevector unit vector is written as
$
\hat{\boldsymbol{q}}
=
\sin\theta_{\boldsymbol{q}}
\left(
\cos\phi_{\boldsymbol{q_\parallel}}\hat{\boldsymbol{x}}
+
\sin\phi_{\boldsymbol{q_\parallel}}\hat{\boldsymbol{y}}
\right)
+
\cos\theta_{\boldsymbol{q}}\hat{\boldsymbol{z}}
\approx
\hat{\boldsymbol{z}},
$
and 
$
\sin\theta_{\boldsymbol{q}} = \frac{q_\parallel}{\sqrt{q_\parallel^2 + q_0^2}}.
$
Substituting the complex amplitude under Lorenz gauge derived previously, each component of the complex amplitude takes the form
\begin{equation} 
\color{G}
\begin{aligned}
\tilde{\mathcal{A}}^{\parallel}_{\ell}(\boldsymbol{q}_\parallel)
&\equiv 
\tilde{\mathcal{A}}^{\parallel}_{\ell, p=0}(\boldsymbol{q}_\parallel)
\approx 
\tilde{F}_{\left|\ell\right|}(q_\parallel)
e^{i\ell \phi_{\boldsymbol{q}_\parallel}} \\
\tilde{\mathcal{A}}^{z}_{\sigma, \ell}(\boldsymbol{q}_\parallel)
&\equiv 
\tilde{\mathcal{A}}^{z}_{\sigma, \ell, p=0}(\boldsymbol{q}_\parallel)
\approx
-\frac{\sin\theta_{\boldsymbol{q}}}{\sqrt{2}}
\tilde{F}_{\left|\ell\right|}(q_\parallel)
e^{i(\sigma+\ell)\phi_{\boldsymbol{q}_\parallel}},
\end{aligned}
\label{AinRinQ} 
\end{equation}

\add{where $\tilde{F}_{\left|\ell\right|}(q_\parallel)=\tilde{F}_{\left|\ell\right|, p=0}(q_\parallel)$.} \add{This expression shows that the SOI is encoded in the phase of the longitudinal component through the combined dependence $(\sigma+\ell)$, corresponding to the total angular momentum (TAM). This behavior difference is illustrated in Figure~\ref{fig:A_C_transeverse} and Figure~\ref{fig:A_C_longitudinal}.} However, this phase structure is generally difficult to detext experimentally, as it is averaged out in intensity-based measurements. To access this TAM-dependent phase, interference between two LG beams can be employed, forming vector vortex beams (VVBs), which will be discussed in the following section.

\begin{figure}[H]
\centering
\includegraphics[width=0.95\linewidth]{A_C_transeverse.pdf}
\caption{\add{Squared magnitude (intensity) and real part distribution of $\tilde{\mathcal{A}}^{\parallel}_{\ell}(\boldsymbol{q}_\parallel)$. The scale bar indicates $0.1$ times the light-cone radius for WSe$_2$ monolayers, $Q_c = 8.6~\mu\mathrm{m}^{-1}$.}}
\label{fig:A_C_transeverse}
\end{figure}

% \add{In Figure~\ref{fig:A_C_longitudinal}, the case $\sigma=0$ is not included 
% because the azimuthal phase factor $e^{i\sigma\phi_{\boldsymbol{q}_\parallel}}$ in Eq.~(\ref{sigma_limit}) 
% is derived under the assumption of circular polarization (so does Eq.~(\ref{AinRinC})), 
% where $\sigma=\pm1$. Therefore, $\sigma=0$ (linear polarization) does not satisfy in this formulation.}

\begin{figure}[H]
\centering
\includegraphics[width=0.95\linewidth]{A_C_longitudinal.pdf}
\caption{\add{Squared magnitude (intensity) and real part distribution of $\tilde{\mathcal{A}}^{z}_{\sigma, \ell}(\boldsymbol{q}_\parallel)$. The scale bar indicates $0.1$ times the light-cone radius for WSe$_2$ monolayers, $Q_c = 8.6~\mu\mathrm{m}^{-1}$.}}
\label{fig:A_C_longitudinal}
\end{figure}

%
%
%
%
%
\subsection{\add{Longitudinal Component in Vector Vortex Beam (VVB)}}

\add{Subject to decoherence,} the information of TAM encoded in the phase \add{is likely} lost when calculating optical transition rates\add{, as pointed out by Oscar Javier Gomez Sanchez \textit{et al.} in Ref.~\cite{sanchez2024enhanced}}. To retain and control this information, they introduced vector vortex beams (VVBs), which are constructed by superimposing two twisted light beams \add{, which in general, these beams carry angular momentaums of the same magnitude but opposite signs. Thus, VVB carries vetoral polarization that varies across the transverse plane. In contrast, scalar vortex beam, single twisted light beams, carries uniform polarization.} The VVB state is expressed as a superposition through
\begin{equation}
\color{G}
\boldsymbol{\mathcal{A}}_{\sigma, \ell, \sigma', \ell'}(\boldsymbol{q}_\parallel;\alpha,\beta)
=
\cos{\left(\frac{\beta}{2}\right)}
\boldsymbol{\mathcal{A}}_{\sigma, \ell}(\boldsymbol{q}_\parallel)
+
e^{i \alpha}
\sin{\left(\frac{\beta}{2}\right)}
\boldsymbol{\mathcal{A}}_{\sigma', \ell'}(\boldsymbol{q}_\parallel).
\end{equation}
The parameter $\alpha$ controls the relative phase between the two beams, while $\beta$ determines their relative weight in the superposition.  \add{Throughout this thesis}, the radial mode index is set to $p=0$, as the primary focus remains on the interplay between SAM and OAM. $\sigma'$ and $\ell'$ denote the second LG beam that carrying a different combination of angular momentum . 

\vspace{0.25ex}
\noindent{\add{\textbf{\,Transverse Components of VVBs}}}\\[0.25ex]
\add{Following Eq.~(\ref{AinCinQ}), a VVB state also contains both transverse and longitudinal components, which can be written as}
\begin{equation}
\color{G}
\boldsymbol{\mathcal{A}}_{\sigma, \ell, \sigma', \ell'}(\boldsymbol{q}_\parallel;\alpha,\beta)
=
\boldsymbol{\mathcal{A}}_{\sigma, \ell, \sigma', \ell'}^\parallel(\boldsymbol{q}_\parallel;\alpha,\beta)
+
\boldsymbol{\mathcal{A}}_{\sigma, \ell, \sigma', \ell'}^z(\boldsymbol{q}_\parallel;\alpha,\beta),
\end{equation}
\add{where the polarization in the transverse component is no longer a fixed unit vector, but becomes$\boldsymbol{q}_\parallel$ dependent as a consequence of the superposition. Therefore, it cannot be factored out as a constant vector in the angular spectrum. The transverse component is given by}
\begin{equation}
\color{G}
\begin{aligned}
&\boldsymbol{\mathcal A}^{\parallel}_{\sigma,\ell,\sigma',\ell'}(\boldsymbol q_\parallel;\alpha,\beta) 
= \mathcal{A}^{x}_{\sigma,\ell,\sigma',\ell'}(\boldsymbol q_\parallel;\alpha,\beta) \hat{\boldsymbol{x}}+\mathcal{A}^{y}_{\sigma,\ell,\sigma',\ell'}(\boldsymbol q_\parallel;\alpha,\beta) \hat{\boldsymbol{y}} \\
&= \frac{1}{\sqrt{2}}\exp\left[ \tfrac{i\left(\ell'+\ell\right)\phi_{\boldsymbol q_\parallel}+i\alpha}{2} \right]  \Bigg\{
\Bigg[
\cos\!\left( \tfrac{\Delta\ell\phi_{\boldsymbol q_\parallel}+\alpha}{2} \right) 
\left( \cos\tfrac{\beta}{2}\tilde F_{|\ell|}(q_\parallel) + \sin\tfrac{\beta}{2}\tilde F_{|\ell'|}(q_\parallel) \right) 
- i\sin\!\left( \tfrac{\Delta\ell\phi_{\boldsymbol q_\parallel}+\alpha}{2} \right) \\
& \left( \cos\tfrac{\beta}{2}\tilde F_{|\ell|}(q_\parallel) - \sin\tfrac{\beta}{2}\tilde F_{|\ell'|}(q_\parallel) \right)
\Bigg] \hat{\boldsymbol x}
+\Bigg[
i\cos\!\left( \tfrac{\Delta\ell\phi_{\boldsymbol q_\parallel}+\alpha}{2} \right)
\left( \sigma\cos\tfrac{\beta}{2}\tilde F_{|\ell|}(q_\parallel) + \sigma'\sin\tfrac{\beta}{2}\tilde F_{|\ell'|}(q_\parallel) \right) \\
& + \sin\!\left( \tfrac{\Delta\ell\phi_{\boldsymbol q_\parallel}+\alpha}{2} \right) 
\left( \sigma\cos\tfrac{\beta}{2}\tilde F_{|\ell|}(q_\parallel) - \sigma'\sin\tfrac{\beta}{2}\tilde F_{|\ell'|}(q_\parallel) \right)
\Bigg] \hat{\boldsymbol y} \Bigg\},
\end{aligned}
\label{VVB_transverse}
\end{equation}
\add{where}
\[
\add{\Delta\ell=\ell'-\ell}
\]
\add{denotes the difference in topological charge between the two constituent beams. Although Eq.~(\ref{VVB_transverse}) appears complicated, its structure can be understood in terms of three distinct contributions: the radial structure, the azimuthal structure, and the global phase.}

\vspace{0.25ex}
\noindent\quad{\add{\textbf{Radial structure}}}~\\*[0.25ex]
\add{The four brackets containing the radial functions $\tilde F_{|\ell|}(q_\parallel)$ determine the radial profile of the VVB. Each twisted light component contributes its own radial distribution, weighted by the angle $\beta$. In particular, when two beams with the same radial size ($|\ell|=|\ell'|$) are equally superposed ($\beta=\pi/2$), the relative sign of the spin $\sigma$ and $\sigma'$ determines the overall polarization state. Specifically, beams with the same spin ($\sigma=\sigma'$) give rise to circular polarization, whereas opposite spin ($\sigma=-\sigma'$) produce linear polarization, as shown below}
\begin{equation}
\color{G}
\begin{aligned}
&\boldsymbol{\mathcal{A}}^{\parallel}_{\sigma, \ell, \sigma', \ell'}\left(\boldsymbol{q}_\parallel; \alpha, \beta = \frac{\pi}{2},\text{with}|\ell|=|\ell'|\right)
= \tilde{F}_{|\ell|}(q_\parallel) 
e^{\frac{i}{2} \left[ (\ell'+\ell)\phi_{\boldsymbol{q}_\parallel}+\alpha \right]} \\
&\times
\left\{
\cos\!\left(
\frac{\Delta\ell\phi_{\boldsymbol q_\parallel}+\alpha}{2}
\right)
\hat{\boldsymbol{x}}
+
\frac{1}{2}
\left[
(\sigma-\sigma')
\sin\!\left(
\frac{\Delta\ell\phi_{\boldsymbol q_\parallel}+\alpha}{2}
\right)
+
(\sigma+\sigma')
i\cos\!\left(
\frac{\Delta\ell\phi_{\boldsymbol q_\parallel}+\alpha}{2}
\right)
\right]
\hat{\boldsymbol y}
\right\},
\end{aligned}
\label{eq:VVB_special_case}
\end{equation}

\vspace{0.25ex}
\noindent\quad{\add{\textbf{Azimuthal structure}}}~\\*[0.25ex]
\add{The trigonometric functions determine the azimuthal variation of the polarization state. For different values of the relative phase $\alpha$, the polarization pattern is uniformly shifted while preserving its overall structure. Within a fixed VVB state, the OAM difference $\Delta\ell$ introduces a periodic spatial modulation along the azimuthal angle. Both behavior can be observed in Fig.~\ref{fig:VVB(transverse_10-10)} ($\Delta\ell=0$) and Fig.~\ref{fig:VVB(transverse_1-1-11)} ($\Delta\ell=2$). In the former case, the polarization remains identical for all azimuthal angles, whereas in the latter case it rotates with the azimuthal angle and repeats periodically. In general, the azimuthal periodicity of the polarization structure is given by}
\[
\add{\text{periodicity}=\frac{2\pi}{\Delta\ell}.}
\]
\add{Therefore, the case $\Delta\ell=2$ shown in Fig.~\ref{fig:VVB(transverse_1-1-11)} exhibits a periodicity of $\pi$. More generally, increasing $|\Delta\ell|$ leads to higher-order polarization textures around the beam axis and reduces the angular period of the polarization pattern.}

\vspace{0.25ex}
\noindent\quad{\add{\textbf{Global Phase}}}~\\*[0.25ex]
\add{Besides the radial and azimuthal structures, Eq.~(\ref{VVB_transverse}) contains an overall phase factor}
\[
\add{{\rm exp}\left[ \tfrac{i\left(\ell'+\ell\right)\phi_{\boldsymbol q_\parallel}+i\alpha}{2} \right]}
\]
\add{which carries both $\alpha$ and the total OAM $\ell+\ell'$ of the two beams. This phase factor does not affect the local polarization state, but contributes to the overall phase winding of the beam. }

\begin{figure}[H]
\centering
\includegraphics[width=1.0\linewidth, trim=0cm 0cm 0cm 0cm, clip]{VVB(transverse_10-10).pdf}
\caption{\add{Real and imaginary parts of $\mathcal{A}^{x}_{1,0,-1,0}(\boldsymbol q_\parallel;\alpha,\beta)$ and $\mathcal{A}^{y}_{1,0,-1,0}(\boldsymbol q_\parallel;\alpha,\beta)$, together with the squared magnitude and polarization distribution of $\boldsymbol{\mathcal A}^{\parallel}_{1,0,-1,0}(\boldsymbol q_\parallel;\alpha,\beta)$.}}
\label{fig:VVB(transverse_10-10)}
\end{figure}

\vspace{0.25ex}
\noindent\quad{\add{\textbf{Intensity}}}~\\*[0.25ex]
\add{Finally, the intensity distribution of the transverse component is given by the squared magnitude}
\begin{equation}
\color{G} 
\begin{aligned}
\left|
\boldsymbol{\mathcal{A}}_{\sigma, \ell, \sigma', \ell'}^\parallel(\boldsymbol{q}_\parallel;\alpha,\beta)
\right|^2
&=
\cos^2\!\left(\frac{\beta}{2}\right)
\left|\tilde{F}_{\left|\ell\right|}(q_\parallel)\right|^2
+
\sin^2\!\left(\frac{\beta}{2}\right)
\left|\tilde{F}_{\left|\ell'\right|}(q_\parallel)\right|^2 \\ 
&\quad+
\frac{1+\sigma\sigma'}{2}\sin\beta\,
\tilde{F}_{\left|\ell\right|}(q_\parallel)
\tilde{F}^{*}_{\left|\ell'\right|}(q_\parallel)
\cos\!\left(\Delta\ell\,\phi_{\boldsymbol{q}_\parallel}+\alpha\right).
\end{aligned}
\label{VVBtrans} 
\end{equation}
\add{For opposite spin ($\sigma=-\sigma'$), the interference contribution in the intensity expression Eq.~(\ref{VVBtrans}) vanishes because $1+\sigma\sigma'=0$. Physically, this reflects the fact that only components with identical polarization can interfere coherently. Consequently, the interference term proportional to $\cos(\Delta\ell\,\phi_{\boldsymbol{q}_\parallel}+\alpha)$ is absent, and the total intensity becomes independent of the relative phase $\alpha$. Furthermore, when the two constituent beams share the same radial profile ($|\ell|=|\ell'|$), their radial functions become identical, and the weighting controlled by the mixing angle $\beta$ no longer affects the intensity distribution. In this case, both $\alpha$ and $\beta$ modify only the polarization structure while leaving the intensity unchanged, as shown below.}
\begin{equation}
\color{G}
\left|
\boldsymbol{\mathcal{A}}_{\sigma, \ell, -\sigma, \ell'}^\parallel(\boldsymbol{q}_\parallel;\alpha,\beta,\text{with}|\ell|=|\ell'|)
\right|^2
=
\left|\tilde{F}_{\left|\ell\right|, p}(q_\parallel)\right|^2 .
\end{equation}

\begin{figure}[H] 
\centering
\includegraphics[width=1.0\linewidth, trim=0cm 0cm 0cm 0cm, clip]{VVB(transverse_1-1-11).pdf} 
\caption{\add{Real and imaginary parts of $\mathcal{A}^{x}_{1,-1,-1,1}(\boldsymbol q_\parallel;\alpha,\beta)$ and $\mathcal{A}^{y}_{1,-1,-1,1}(\boldsymbol q_\parallel;\alpha,\beta)$, together with the squared magnitude and polarization distribution of $\boldsymbol{\mathcal A}^{\parallel}_{1,-1,-1,1}(\boldsymbol q_\parallel;\alpha,\beta)$.}}
\label{fig:VVB(transverse_1-1-11)}
\end{figure}
\vspace{0.25ex}
\noindent{\add{\textbf{\,Transverse Components of VVBs}}}\\[0.25ex]
\add{The longitudinal component is given by}
\begin{equation}
\color{G}
\begin{aligned}
&\boldsymbol{\mathcal{A}}_{\sigma, \ell, \sigma', \ell'}^z(\boldsymbol{q}_\parallel;\alpha,\beta) 
= \mathcal{A}^{z}_{\sigma,\ell,\sigma',\ell'}(\boldsymbol q_\parallel;\alpha,\beta)\,\hat{\boldsymbol{z}}\\
&= \frac{-\sin\!\theta_{\boldsymbol{q}}}{\sqrt{2}}\exp\left[ \tfrac{i\left(\ell'+\sigma'+\ell+\sigma\right)\phi_{\boldsymbol q_\parallel}+i\alpha}{2} \right]  
\Bigg\{
\Bigg[
\cos\!\left( \tfrac{\Delta J \phi_{\boldsymbol q_\parallel}+\alpha}{2} \right) 
\left( \cos\tfrac{\beta}{2}\tilde F_{|\ell|}(q_\parallel) + \sin\tfrac{\beta}{2}\tilde F_{|\ell'|}(q_\parallel) \right) \\
& \hspace{14em}
- i\sin\!\left( \tfrac{\Delta J \phi_{\boldsymbol q_\parallel}+\alpha}{2} \right)
\left( \cos\tfrac{\beta}{2}\tilde F_{|\ell|}(q_\parallel) - \sin\tfrac{\beta}{2}\tilde F_{|\ell'|}(q_\parallel) \right)
\Bigg]
\Bigg\}\hat{\boldsymbol z}
\end{aligned}
\label{VVB_longitudinal}
\end{equation}
\noindent
\add{where}
\[
\add{\Delta J =(\ell'+\sigma')-(\ell+\sigma)}
\]
\add{denotes the difference in TAM between the two constituent beams, which reflecting its origin from the longitudinal field in Eq.~(\ref{AinRinQ}). This expression has a structure similar to Eq.~(\ref{VVB_transverse}), and can likewise be decomposed into radial, azimuthal, and global phase contributions. The key difference is that the OAM index is replaced by the TAM indices. As a result, the azimuthal dependence exhibits a periodicity of}
\[
\add{\text{periodicity} = \frac{2\pi}{\Delta J}.}
\]
\noindent
\add{In addition, the exponential prefactor carries the global phase associated with the total angular momentum of the superposed state, thereby encoding the combined spin and orbital contributions.}
\noindent
\add{In a similar way, the intensity distribution of the longitudinal component is given by the squared magnitude of} % 修正為 In a similar way... of
\begin{equation}
\color{G} 
\begin{aligned}
\left|
\boldsymbol{\mathcal{A}}_{\sigma, \ell, \sigma', \ell'}^z(\boldsymbol{q}_\parallel;\alpha,\beta)
\right|^2
&=
\frac{\sin^2\theta_{\boldsymbol{q}}}{2}
\Bigg[
\cos^2\!\left(\frac{\beta}{2}\right)
\left|\tilde{F}_{\left|\ell\right|}(q_\parallel)\right|^2
+
\sin^2\!\left(\frac{\beta}{2}\right)
\left|\tilde{F}_{\left|\ell'\right|}(q_\parallel)\right|^2 \\ 
&\qquad+
\sin\beta\,
\tilde{F}_{\left|\ell\right|}(q_\parallel)\tilde{F}^{*}_{\left|\ell'\right|}(q_\parallel)
\cos\!\left(\Delta J\,\phi_{\boldsymbol{q}_\parallel}+\alpha\right)
\Bigg],
\end{aligned}
\label{VVBlongi} 
\end{equation}
\add{Since the longitudinal component does not involve a polarization superposition in the same way as the transverse component, there is no constraint from polarization coherence that suppresses interference. As a result, the interference term proportional to $\cos(\Delta J \,\phi_{\boldsymbol{q}_\parallel}+\alpha)$ is preserved, as explicitly shown in Eq.~(\ref{VVBlongi}). Similarly, when the two beams share the same radial profile ($|\ell|=|\ell'|$), we can obtain a cleaner formulation dedicated to the azimuthal variations.} 
\begin{equation}
\color{G}
\begin{aligned}
\left|
\boldsymbol{\mathcal{A}}_{\sigma, \ell, -\sigma, \ell}^z(\boldsymbol{q}_\parallel;\alpha,\beta)
\right|^2
&=
\frac{\sin^2\theta_{\boldsymbol{q}}}{2}
\left|\tilde{F}_{\left|\ell\right|, p}(q_\parallel)\right|^2
\left[
1+\sin\beta \,
\cos\!\left(\Delta J\,\phi_{\boldsymbol{q}_\parallel}+\alpha\right)
\right].
\end{aligned}
\label{VVBlongi-simple}
\end{equation}

In this case, $\Delta J=-2\sigma$. The transverse component therefore exhibits an isotropic squared magnitude in angular spectrum and does not retain information about the SOI. By contrast, the longitudinal component reflects the TAM difference between the two LG beams through the interference term. Here, both beams share the same radial factor $\tilde{F}_{\left|\ell\right|, p}(q_\parallel)$ for $|\ell'|=|\ell|$. As a result, the common radial factor can be factored out, allowing us to focus directly on the angular interference. 

For a simple case $\sigma'=-\sigma=-1$ and $\ell'=\ell=0$, Figure~\ref{fig:VVBs} shows that the squared magnitude of the longitudinal component exhibits TAM dependent azimuthal modulation as $\alpha$ and $\beta$ are varied. In contrast, the transverse components retain the same intensity distribution over the Poincar\'e sphere, while the variation appears only in the polarization texture. Note that the factor $\sin^2\theta_{\boldsymbol{q}}$ in Eq.~(\ref{VVBlongi-simple}) vanishes near $\boldsymbol{q}_\parallel\approx0$, producing a central null. As a result, the longitudinal component preserves a donut-shaped profile even when $\ell=\ell'=0$ for both LG beams.


Motivated by this behavior, one may instead choose $\sigma'=-\sigma=-1$ and $\ell'=\ell=-1$, for which $\Delta J=(\ell'+\sigma')-(\ell+\sigma)=0$. In this case, the longitudinal component becomes azimuthally uniform and can exhibit fully constructive or destructive interference over the entire angular spectrum when $(\alpha,\beta)=\left(0,\frac{\pi}{2}\right)$ or $\left(\pi,\frac{\pi}{2}\right)$, as shown in Figure~\ref{fig:VVBs_lm1}.















